\documentclass[12pt]{article}
\usepackage[T1]{fontenc}
\usepackage[utf8]{inputenc}
\usepackage[polish]{babel}
\usepackage{lmodern}
\usepackage[a4paper,margin=2.4cm]{geometry}
\usepackage{xcolor}
\usepackage{graphicx}
\usepackage{amsmath,amssymb}
\usepackage{booktabs}
\usepackage{array}
\usepackage{enumitem}
\usepackage{hyperref}
\usepackage{tcolorbox}
\usepackage{fancyhdr}
\definecolor{WEIIPrimary}{RGB}{0,130,150}
\definecolor{WEIIPrimaryDark}{RGB}{0,86,100}
\definecolor{WEIIPrimaryLight}{RGB}{229,246,248}
\definecolor{PRzNavy}{RGB}{31,63,142}
\definecolor{RLTextDark}{RGB}{34,34,34}
\definecolor{RLGray}{RGB}{90,90,90}
\definecolor{RLLightGray}{RGB}{245,246,248}
\definecolor{RLWarning}{RGB}{184,86,0}
\definecolor{RLExample}{RGB}{52,112,64}
\definecolor{RLCodeBg}{RGB}{248,248,248}

\usepackage{listings}
\usepackage{xcolor}

\lstdefinestyle{rlpython}{%
    language=Python,
    basicstyle=\ttfamily\scriptsize,
    keywordstyle=\color{WEIIPrimaryDark}\bfseries,
    commentstyle=\color{RLGray}\itshape,
    stringstyle=\color{RLExample},
    numbers=left,
    numberstyle=\tiny\color{RLGray},
    numbersep=6pt,
    showstringspaces=false,
    breaklines=true,
    tabsize=4,
    frame=single,
    rulecolor=\color{WEIIPrimary!35},
    backgroundcolor=\color{RLCodeBg},
    captionpos=b
}

\lstset{style=rlpython}


\hypersetup{colorlinks=true,linkcolor=WEIIPrimaryDark,urlcolor=WEIIPrimaryDark,citecolor=WEIIPrimaryDark}
\setlength{\headheight}{15pt}
\pagestyle{fancy}
\fancyhf{}
\lhead{Uczenie się ze wzmocnieniem}
\rhead{WEiI PRz}
\cfoot{\thepage}

\newtcolorbox{taskbox}[1]{colback=WEIIPrimaryLight,colframe=WEIIPrimary,title=\textbf{#1}}
\newtcolorbox{notebox}[1]{colback=RLLightGray,colframe=RLGray,title=\textbf{#1}}
\newtcolorbox{warninghandout}[1]{colback=RLWarning!8,colframe=RLWarning,title=\textbf{#1}}

\newcommand{\handouttitle}[3]{%
\begin{center}
{\Large\bfseries #1}\\[0.2cm]
{\large #2}\\[0.15cm]
{\small #3}
\end{center}
\vspace{0.4cm}
\hrule
\vspace{0.6cm}
}

\newcommand{\EnableLabNumbering}{%
  \renewcommand{\thesection}{L\arabic{section}}%
  \renewcommand{\thesubsection}{\thesection.\arabic{subsection}}%
  \renewcommand{\thesubsubsection}{\thesubsection.\arabic{subsubsection}}%
}


\title{Raport projektowy\\Uczenie się ze wzmocnieniem\\[0.35em]\large Propozycja tematu -- etap E1}
\author{\textbf{Alan Kyciński}}
\date{Rok akademicki 2026/2027}

\begin{document}
\maketitle
\tableofcontents
\newpage


\section{Informacje podstawowe}

\begin{center}
\renewcommand{\arraystretch}{1.22}
\begin{tabular}{p{0.34\textwidth}p{0.56\textwidth}}
\toprule
Element & Informacja \\
\midrule
Imię i nazwisko & \textbf{Alan Kyciński} \\
Numer albumu & \textbf{179545} \\
Kierunek/specjalność & Informatyka / Sztuczna inteligencja \\
Tytuł projektu & Optymalizacja procesu kontroli jakości produktów przemysłowych z wykorzystaniem algorytmu Q-learning \\
Link do repozytorium & {\small\url{https://github.com/KciamiK/DefectDecision-RL}} \\
Data oddania & \textbf{6 listopada 2026 r.} \\
\bottomrule
\end{tabular}
\end{center}

\section{Temat i cel projektu}

\subsection{Opis problemu}

Projekt \textbf{DefectDecision-RL} dotyczy optymalizacji procesu przemysłowej kontroli jakości z wykorzystaniem uczenia się ze wzmocnieniem. Kontrolowanym wyrobem jest \textbf{prostokątna płytka aluminiowa z czterema otworami montażowymi}. Kolejne płytki mogą być prawidłowe, posiadać dopuszczalne ślady obróbki lub uszkodzenia przekraczające ustalone wymagania jakościowe. Zadaniem systemu jest podejmowanie decyzji dotyczących dalszego postępowania z kontrolowanymi produktami.

Rozważane są \textbf{cztery grupy nieprawidłowości}: (1) rysy powierzchniowe, (2) pęknięcia, (3) nieprawidłowe otwory montażowe oraz (4) plamy i ślady korozji. Na jednej płytce może występować kilka nieprawidłowości jednocześnie. Nie każda z nich oznacza jednak wadliwy produkt. Drobne zarysowania, niewielkie przebarwienia czy odchylenia położenia otworów mogą mieścić się w przyjętych tolerancjach. Natomiast rzeczywiste pęknięcia, brak otworu lub przekroczenie dopuszczalnych odchyleń powodują uznanie produktu za wadliwy. Wszystkie kryteria jakościowe zostaną określone na potrzeby symulacji i nie będą stanowiły rzeczywistych norm przemysłowych.

W środowisku generowane będą \textbf{syntetyczne obrazy 2D płytek} uwzględniające różne rodzaje i stopnie nasilenia uszkodzeń, a także zmienność oświetlenia, faktury powierzchni i szum obrazu. Każdy produkt zostanie poddany wstępnemu skanowaniu, którego wynikiem będzie oszacowanie ryzyka wystąpienia niedopuszczalnej wady. Na podstawie otrzymanych informacji agent wybierze jedną z trzech akcji: \texttt{ACCEPT} (akceptacja), \texttt{REJECT} (odrzucenie) lub \texttt{INSPECT} (dodatkowa kontrola). Szczegółowe badanie dostarczy dokładniejszych informacji, jednak jego wykonanie będzie związane z ograniczonym budżetem inspekcji.

\textbf{Głównym celem projektu jest opracowanie i zbadanie strategii decyzyjnej opartej na algorytmie Q-learning.} Agent będzie uczył się podejmowania decyzji w warunkach niepewności, uwzględniając skuteczność kontroli oraz konsekwencje wykorzystania dostępnych zasobów. Opracowane rozwiązanie zostanie poddane eksperymentom pozwalającym ocenić jego zdolność do ograniczania liczby błędnych decyzji i kosztów inspekcji. Uzyskane wyniki zostaną porównane z rezultatami prostych strategii bazowych.

\subsection{Cel agenta}

Agent ma nauczyć się \textbf{dobierać decyzję do niepewnego wyniku badania, liczby pozostałych płytek i budżetu szczegółowych kontroli}. Celem jest maksymalizacja sumy nagród w całej partii, a nie uzyskanie jak największej liczby pozytywnych wskazań skanera. System powinien:
\begin{itemize}[nosep]
    \item skutecznie identyfikować płytki przekraczające tolerancje jakościowe;
    \item nie odrzucać bez potrzeby płytek z dopuszczalnymi śladami produkcji;
    \item ograniczać koszt dodatkowych badań i zachowywać budżet na kolejne przypadki;
    \item podejmować rozsądne decyzje również przy niejednoznacznych pomiarach.
\end{itemize}

Wyniki oceny będą obejmować średnią nagrodę na partię, udział przeoczonych wad, udział błędnych odrzuceń i liczbę wykonanych badań szczegółowych. Agent Q-learning zostanie porównany z agentem losowym i heurystyką opartą na progach ryzyka. Interfejs w Pygame pokaże produkt, wynik skanu, decyzję agenta, stan budżetu oraz przebieg uczenia i porównania.

\subsection{Uzasadnienie zastosowania uczenia się ze wzmocnieniem}

Problem jest \textbf{sekwencyjny}: w ramach jednej partii agent kontroluje kolejne produkty, a po zleceniu badania szczegółowego może podjąć następną decyzję dotyczącą tej samej płytki. Decyzja o inspekcji wpływa na przyszły stan środowiska, ponieważ zmniejsza budżet pozostały dla kolejnych wyrobów. Decyzje akceptacji i odrzucenia są ostateczne, a ocena skanera bywa niepewna.

Występuje zatem \textbf{konflikt celów}: ograniczenie błędów jakościowych wymaga czasem dodatkowej analizy, ale zasoby na taką analizę są ograniczone. Prosty klasyfikator działający niezależnie dla każdej płytki nie uwzględnia kosztu utraconej możliwości kontroli w przyszłości. Q-learning będzie uczył się strategii maksymalizującej \textbf{długoterminowy zwrot} poprzez wielokrotne epizody i informację zwrotną w postaci nagród i kar. Porównanie z heurystyką pozwoli sprawdzić, czy uczenie jest uzasadnione również praktycznie.

\section{Formalizacja problemu RL/MDP}

\subsection{Agent i środowisko}

\textbf{Agent}: dyskretny kontroler inspekcji, który wybiera działanie dla aktualnej płytki na podstawie obserwowanych wyników badań i pozostałego budżetu.

\textbf{Środowisko}: autorski symulator generujący płytki 2D, ich ukrytą zgodność z wymaganiami, wyniki skanowania i nagrody. Generator przechowuje rzeczywiste parametry defektów (np. geometrię rysy i otworów), a reguły tolerancji wyznaczają etykietę \(y\in\{0,1\}\), gdzie \(0\) oznacza płytkę zgodną z wymaganiami, a \(1\) -- wadliwą. \textbf{Etykieta \(y\) nie jest przekazywana agentowi przed końcową decyzją}. Używa jej środowisko do naliczenia nagrody.

Każdy krok realizuje interfejs \texttt{reset()} oraz \texttt{step(action)}. Po \texttt{ACCEPT} albo \texttt{REJECT} środowisko rozlicza decyzję i przechodzi do następnej płytki. Po \texttt{INSPECT} zużywa jedną kontrolę, udostępnia dodatkowy, niedoskonały wynik skanera i pozostaje przy tym samym produkcie. Skanery nie wykonują za agenta decyzji końcowej.

\subsection{Stan lub obserwacja}

Pełny stan środowiska zawiera m.in. ukryte parametry defektów i faktyczną zgodność płytki. Agent otrzymuje wyłącznie \textbf{dyskretną obserwację}:
\begin{equation}
    s_t=(n_t,b_t,w_t,z_t),
\end{equation}
\noindent gdzie:
\begin{center}
\renewcommand{\arraystretch}{1.2}
\begin{tabular}{p{0.29\textwidth}p{0.61\textwidth}}
\toprule
Składowa & Znaczenie \\
\midrule
\(n_t\in\{1,\ldots,10\}\) & Liczba płytek do rozstrzygnięcia, wliczając bieżącą. \\
\(b_t\in\{0,\ldots,4\}\) & Pozostała liczba badań szczegółowych w partii. \\
\(w_t\in\{L,M,H\}\) & Niska, średnia lub wysoka kategoria ryzyka po skanie podstawowym. \\
\(z_t\in\{\varnothing,-,+\}\) & Brak dodatkowego badania albo jego wynik: brak wskazania wady / wskazanie wady. \\
\bottomrule
\end{tabular}
\end{center}

Przykład \(s=(7,3,M,\varnothing)\) oznacza siedem płytek pozostałych do decyzji, trzy dodatkowe kontrole i średnie ryzyko dla płytki dotychczas badanej tylko podstawowo. Wariant \(z\neq\varnothing\) sygnalizuje, że dodatkowy skan już wykonano. Pełne wyniki czterech analizowanych cech będą widoczne w panelu informacyjnym, natomiast \textbf{do tablicy Q trafi ich zwięzłe podsumowanie}. Utrzymuje to niewielką przestrzeń stanów.

W formalnym ujęciu jakość produktu jest częściowo obserwowalna (problem zbliżony do POMDP). Używana obserwacja \(s_t\) będzie \textbf{przybliżeniem stanu decyzyjnego}, opartym na skategoryzowanych wynikach badań i zasobach. Przy stałych parametrach symulacji posłuży do tablicowego Q-learningu, a ograniczenie wynikające z utraty szczegółów pomiaru zostanie omówione w raporcie. Tabela przechowuje maksymalnie \(10\cdot5\cdot3\cdot3\cdot3=1350\) wartości \(Q(s,a)\), z których część odpowiada niedozwolonym lub nieosiągalnym przypadkom.

\subsection{Akcje}

\begin{center}
\renewcommand{\arraystretch}{1.22}
\begin{tabular}{p{0.28\textwidth}p{0.62\textwidth}}
\toprule
Akcja & Działanie i następstwa \\
\midrule
\(a=0\), \texttt{ACCEPT} & Przyjęcie płytki; nagroda za zgodną z prawdą decyzję albo kara za przeoczenie wady; przejście do następnej płytki. \\
\(a=1\), \texttt{REJECT} & Odrzucenie płytki; nagroda za wykrycie wady albo kara za odrzucenie prawidłowego wyrobu; przejście do następnej płytki. \\
\(a=2\), \texttt{INSPECT} & Dodatkowy skan bieżącej płytki, zużycie jednostki budżetu, nowy wynik badania i ponowny wybór decyzji. \\
\bottomrule
\end{tabular}
\end{center}

Akcja \texttt{INSPECT} jest \textbf{niedostępna}, gdy \(b_t=0\) lub \(z_t\neq\varnothing\), czyli płytka została już sprawdzona dodatkowo. Nie ma możliwości wykonania więcej niż jednego dodatkowego badania tej samej płytki. Wybór akcji i aktualizacja Q będą uwzględniać wyłącznie legalne działania; pozostałe akcje są zawsze dostępne.

\subsection{Funkcja nagrody}

Środowisko oblicza nagrodę na podstawie wybranej akcji i \textbf{rzeczywistej zgodności z tolerancjami}, a nie na podstawie samego wyniku skanera. Przyjmujemy następujące \textbf{wstępne} wartości:
\begin{equation}
 r_t=\begin{cases}
 +2, & a_t=\mathrm{ACCEPT}\ \land\ y_t=0,\\
 -12, & a_t=\mathrm{ACCEPT}\ \land\ y_t=1,\\
 +5, & a_t=\mathrm{REJECT}\ \land\ y_t=1,\\
 -4, & a_t=\mathrm{REJECT}\ \land\ y_t=0,\\
 -1, & a_t=\mathrm{INSPECT}.
 \end{cases}
\end{equation}

Największa kara \(-12\) dotyczy przepuszczenia wadliwej płytki. Odrzucenie płytki mieszczącej się w tolerancji również jest błędem (kara \(-4\)), nawet gdy ma ona widoczne drobne rysy lub przebarwienia. Inspekcja kosztuje \(-1\), dlatego agent powinien ją wybierać tylko wtedy, gdy wartość dodatkowej informacji uzasadnia koszt teraz i stratę części budżetu na przyszłość. Poprawne rozstrzygnięcia przynoszą nagrodę dodatnią.

Nie wprowadzamy dodatkowej kary za każdy krok. Dopuszczalność akcji i ograniczona liczba inspekcji zapobiegają niekończącym się epizodom. \textbf{Wagi nagród będą badane eksperymentalnie}; możliwe niepożądane strategie, np. zbyt częste odrzucanie produktów, zostaną ocenione osobno przez metryki błędów.

\subsection{Warunki zakończenia epizodu}

\textbf{Jeden epizod to jedna partia 10 aluminiowych płytek}. Na początku epizodu losowana jest nowa partia, a budżet szczegółowych badań wynosi \textbf{4}. Epizod kończy się po podjęciu ostatecznej decyzji \texttt{ACCEPT} albo \texttt{REJECT} dla dziesiątej płytki. Wyczerpanie budżetu nie kończy epizodu: dalsze płytki trzeba rozpatrzyć bez dodatkowego skanowania.

W jednym epizodzie występuje \textbf{10 decyzji końcowych i najwyżej 4 dodatkowe skany}, czyli łącznie od 10 do 14 akcji. Za ostatnią decyzję przyznawana jest normalna nagroda, po czym środowisko zwraca znacznik końca epizodu. Liczba działań jest ograniczona i nie ma ryzyka nieskończonej sekwencji.

\subsection{Założenia i ograniczenia środowiska}

\textbf{Produkt i tolerancje.} W projekcie kontrolowany będzie jeden typ wyrobu: płytka z czterema otworami. Generator tworzy obrazy prawidłowych płytek oraz cztery grupy nieprawidłowości o losowym natężeniu. O etykiecie produktu decydują \textbf{zdefiniowane wcześniej reguły jakościowe}, np. maksymalna wielkość rysy, dopuszczalne odchylenie otworu i zakres przebarwienia. Prawdziwe pęknięcie i brak otworu zawsze oznaczają wadę; drobne ślady obróbki mogą być akceptowalne. Tolerancje służą wyłącznie do zbudowania wiarygodnej symulacji, nie do certyfikowania rzeczywistych części. Możliwe będzie współwystępowanie kilku śladów na jednej płytce.

\textbf{Wyniki skanerów.} Skaner podstawowy oceni w uproszczeniu cztery obszary: linie i rysy, nieregularne pęknięcia, geometrię otworów oraz plamy. Będzie działał na mniej szczegółowym obrazie; skaner dodatkowy skorzysta z dokładniejszego pomiaru. Oba będą obarczone szumem i mogą popełniać błędy. Do analizy 2D przewiduje się proste operacje na obrazie (np. OpenCV), bez trenowania detektora sieciowego. Skaner \textbf{nie otrzyma z generatora informacji o prawdziwych wadach}; analizować będzie tylko obraz. Wstępnie dopuszcza się model statystyczny wyników badań do przetestowania samego środowiska RL, zanim powstanie pełny prosty skaner obrazu.

\textbf{Wynik ryzyka.} Każdy skaner będzie przedstawiał wyniki pomiarów cząstkowych; na ich podstawie zostanie wyznaczony zbiorczy wynik ryzyka. Wyniku nie należy przedstawiać jako rzeczywistego \emph{procentowego prawdopodobieństwa} bez kalibracji. Początkowo będzie to wynik umowny 0--100, a dla Q-learningu zostanie zamieniony na kategorie \(L,M,H\). Jeśli pozwolą na to wyniki testów, oszacowanie prawdopodobieństwa wady będzie kalibrowane na oddzielnym zbiorze syntetycznych przykładów. Do agenta nie trafi prawdziwa etykieta defektu ani informacja o przekroczonych tolerancjach.

\textbf{Ograniczenia i realizacja etapowa.} Model pomija fizyczne własności metalu, rzeczywiste kamery, koszty produkcji oraz złożone defekty przemysłowe. \textbf{Priorytetem jest działający, zweryfikowany agent RL}: (1) generator płytek i etykiet oraz testy środowiska, (2) podstawowy i dodatkowy model obserwacji, (3) agenci losowy, heurystyczny i Q-learning, (4) prosty skaner syntetycznych obrazów, wizualizacja Pygame i eksperymenty. Jeśli etap analizy obrazów okaże się zbyt złożony, będzie uproszczony bez zmiany problemu decyzyjnego i formalizacji RL.

\textbf{Plan uczenia i oceny.} Zostanie zastosowany tablicowy Q-learning z eksploracją \(\varepsilon\)-greedy. Wstępnie przyjmuje się \(\alpha=0{,}1\), \(\gamma=0{,}95\), stopniowe zmniejszanie \(\varepsilon\) od \(1{,}0\) do \(0{,}05\) oraz około 20\,000 epizodów treningowych. Parametry są propozycją, nie ustalonym wynikiem badań. W eksperymentach planuje się porównanie z agentem losowym i regułowym, zmianę budżetu oraz trudności skanowania i analizę liczby przeoczonych wad, błędnych odrzuceń, średniej nagrody i wykorzystania badań. Losowania treningowe i testowe będą rozdzielone, a ziarna losowe zapisane, aby wyniki dało się odtworzyć.

\section{Implementacja}

\subsection{Struktura repozytorium}

Należy krótko opisać strukturę repozytorium i najważniejsze pliki projektu.

\begin{verbatim}
repozytorium/
|-- src/
|-- experiments/
|-- report/
|-- results/
`-- README.md
\end{verbatim}

\subsection{Środowisko}

Należy opisać sposób implementacji środowiska, w szczególności funkcje inicjalizacji epizodu oraz wykonania kroku środowiska.

\subsection{Agent bazowy}

Należy opisać agenta losowego, heurystycznego albo inny baseline wykorzystany do porównania z agentem uczącym się.

\subsection{Agent uczący się}

Należy opisać zastosowany algorytm RL lub DRL, jego parametry oraz sposób implementacji. W przypadku Q-tab należy opisać reprezentację tablicy wartości. W przypadku DRL należy opisać wejścia modelu, wyjścia modelu, architekturę oraz sposób uczenia.

\section{Plan eksperymentów}

Należy opisać eksperymenty zaplanowane w celu oceny działania agenta.

\begin{center}
\begin{tabular}{p{0.24\textwidth}p{0.30\textwidth}p{0.30\textwidth}}
\toprule
Eksperyment & Co jest zmieniane? & Co jest mierzone? \\
\midrule
E1 & Do uzupełnienia & Do uzupełnienia \\
E2 & Do uzupełnienia & Do uzupełnienia \\
E3 & Do uzupełnienia & Do uzupełnienia \\
\bottomrule
\end{tabular}
\end{center}

W tej sekcji należy wskazać:
\begin{itemize}
    \item liczbę epizodów uczenia,
    \item liczbę powtórzeń eksperymentu albo uzasadnienie braku powtórzeń,
    \item parametry uczenia,
    \item metryki oceny,
    \item warianty środowiska, map lub funkcji nagrody,
    \item sposób porównania z agentem bazowym.
\end{itemize}

\section{Wyniki}

W tej sekcji należy przedstawić wyniki eksperymentów. Wyniki powinny być pokazane w postaci tabel, wykresów lub czytelnych zestawień.

\begin{center}
\begin{tabular}{p{0.28\textwidth}p{0.20\textwidth}p{0.20\textwidth}p{0.20\textwidth}}
\toprule
Wariant & Średni zwrot & Skuteczność & Długość epizodu \\
\midrule
Agent bazowy & Do uzupełnienia & Do uzupełnienia & Do uzupełnienia \\
Agent RL/DRL & Do uzupełnienia & Do uzupełnienia & Do uzupełnienia \\
\bottomrule
\end{tabular}
\end{center}

\begin{figure}[h]
    \centering
    \fbox{\parbox{0.82\textwidth}{\centering Miejsce na wykres procesu uczenia}}
    \caption{Przykładowy wykres procesu uczenia agenta.}
\end{figure}

\section{Analiza i wnioski}

Należy zinterpretować wyniki. W tej części należy odpowiedzieć m.in. na pytania:
\begin{itemize}
    \item Czy agent uczący się działa lepiej niż agent bazowy?
    \item Jak zmieniało się zachowanie agenta w trakcie uczenia?
    \item Które elementy środowiska były najtrudniejsze?
    \item Czy funkcja nagrody prowadziła do oczekiwanego zachowania?
    \item Czy wystąpiły przypadki niepożądanego zachowania agenta?
    \item Jakie są najważniejsze ograniczenia rozwiązania?
\end{itemize}

\section{Instrukcja uruchomienia}

Należy opisać sposób uruchomienia projektu od czystego repozytorium. Przykładowo:

\begin{verbatim}
git clone https://github.com/...
cd nazwa-repozytorium
python -m venv .venv
.venv\Scripts\activate
pip install -r requirements.txt
python src/main.py
\end{verbatim}

Należy również wskazać, jak odtworzyć eksperymenty oraz gdzie zapisywane są wyniki.

\section{Wsparcie AI}

W tej sekcji należy obowiązkowo opisać wykorzystanie narzędzi sztucznej inteligencji. Jeżeli student nie korzystał z takich narzędzi, należy wpisać jednoznaczną informację, że nie korzystano ze wsparcia AI.

\begin{center}
\begin{tabular}{p{0.13\textwidth}p{0.19\textwidth}p{0.21\textwidth}p{0.27\textwidth}}
\toprule
Data & Narzędzie/model & Cel użycia & Transkrypt / link / zakres \\
\midrule
Do uzupełnienia & Do uzupełnienia & Do uzupełnienia & Do uzupełnienia \\
Do uzupełnienia & Do uzupełnienia & Do uzupełnienia & Do uzupełnienia \\
\bottomrule
\end{tabular}
\end{center}

Należy zapisać transkrypty rozmów z modelami AI wraz z datą oraz nazwą modelu. Jeżeli rozmowa odbywała się poza narzędziami programistycznymi, należy dodatkowo udostępnić całą rozmowę poprzez link, o ile użyte narzędzie umożliwia takie udostępnienie.

Student powinien także wyjaśnić, które elementy wygenerowane lub zasugerowane przez AI zostały samodzielnie sprawdzone, zmodyfikowane albo odrzucone.

\section{Literatura}

W tej sekcji należy wskazać wykorzystane źródła: książki, artykuły naukowe, dokumentację bibliotek, materiały wykładowe oraz inne źródła użyte podczas pracy nad projektem.

\begin{thebibliography}{9}
\bibitem{sutton_barto}
R. S. Sutton, A. G. Barto, \textit{Reinforcement Learning: An Introduction}, 2nd ed., MIT Press, 2018.
\end{thebibliography}

\end{document}
